\questao{94}{corA}{Mecânica}{Quantidade de movimento e energia cinética}{Com conta}
\begin{enunciado}
% gerado por ferramentas/md2tex_enunciado.py a partir de questoes/q094.md (não editar à mão)
Em um reator de fusão nuclear, é necessário produzir plasma a um milhão de graus centígrados para desencadear reações de fusão nuclear controladas. O plasma é formado por um gás ionizado que contém inicialmente deutério $D^{+}$ e trítio $T^{+}$ (isótopos do hidrogênio, com um e dois nêutrons, respectivamente). Nesse plasma, a energia térmica é suficientemente alta para desencadear uma reação de fusão nuclear na qual o núcleo do deutério, partícula $D^{+}$, e o núcleo de trítio, partícula $T^{+}$, colidem para formar uma partícula $\alpha$ e um nêutron \textit{n}. A representação dessa reação de fusão é\par
\[ D^{+} + T^{+} \longrightarrow \alpha + n \]
A partícula $\alpha$ é formada por dois prótons e dois nêutrons. A energia liberada pela fusão é dividida na forma de energia cinética entre a partícula $\alpha$ e o nêutron \textit{n}. Considere as massas do próton e do nêutron iguais, que o momento linear total das partículas reagentes é nulo e também que a energia cinética dos produtos é de 20\,MeV.\par
\fonte{Disponível em: www.scientificamerican.com. Acesso em: 11 nov. 2019 (adaptado).}
Nessa reação de fusão, a energia cinética do nêutron, em megaelétron-volt, é\par
\begin{altlinha}
\alth{A}{4,0\,MeV.}%
\alth{B}{10\,MeV.}%
\alth{C}{13\,MeV.}%
\alth{D}{15\,MeV.}%
\alth{E}{16\,MeV.}
\end{altlinha}

\end{enunciado}

\quer{dos 20\un{MeV} de energia liberados na fusão, quanto fica com o nêutron. (MeV, megaelétron-volt, é só uma unidade de energia, como o joule. Não é preciso converter.)}

\begin{ideia}
A questão fala de fusão nuclear, mas a conta é de \textbf{Mecânica}. A frase que resolve tudo é:
``o momento linear total das partículas reagentes é nulo''.

\textbf{Momento linear} é outro nome da \textbf{quantidade de movimento}: massa vezes velocidade ($Q = m \cdot v$).
Ele se conserva: se era zero antes da reação, continua zero depois. Para a soma dar zero, as duas partículas que saem
precisam ir \textbf{em sentidos opostos, com a mesma quantidade de movimento}.

\textbf{Pense em dois patinadores parados que se empurram:} um adulto pesado e uma criança leve. Os dois recebem o
mesmo empurrão, mas a criança sai muito mais rápida. E quem sai mais rápido leva a maior parte da energia.
Aqui, a partícula $\alpha$ é o adulto e o nêutron é a criança.

\begin{center}
\begin{tikzpicture}[font=\small]
  \node[circle,fill=corA!30,draw=corA,minimum size=11mm,inner sep=0pt] (a) at (-2.0,0) {$\alpha$};
  \node[circle,fill=corB!30,draw=corB,minimum size=5.5mm,inner sep=0pt] (n) at (2.0,0) {$n$};
  \draw[seta,corA] (a.west) -- ++(-1.0,0) node[pos=1,left]{devagar};
  \draw[seta,corB] (n.east) -- ++(4.0,0) node[midway,above]{4 vezes mais rápido};
  \node[below=1mm of a,align=center,cinza] (ma) {massa 4\\(2 prótons + 2 nêutrons)};
  \node[anchor=north,cinza] at (n |- ma.north) {massa 1};
\end{tikzpicture}
\end{center}
\end{ideia}

\passoapasso
\npasso{Compare as massas.}
O enunciado manda considerar a massa do próton igual à do nêutron. Vamos chamar essa massa de \textbf{1 unidade}.
\begin{itemize}
  \item Nêutron: massa 1.
  \item Partícula $\alpha$: 2 prótons + 2 nêutrons = 4 partículas. Massa 4.
\end{itemize}

\npasso{Use a conservação da quantidade de movimento para comparar as velocidades.}
As duas quantidades de movimento são iguais (em sentidos opostos): massa $\times$ velocidade da $\alpha$ = massa $\times$ velocidade do nêutron.
\begin{conta}
$4 \times v_\alpha = 1 \times v_n$
\end{conta}
Para os dois lados ficarem iguais, a velocidade do nêutron tem de ser 4 vezes a da $\alpha$.
O nêutron é 4 vezes mais leve, então sai \textbf{4 vezes mais rápido}.

\npasso{Compare as energias cinéticas.}
Energia cinética é a energia do movimento: $E_c = \dfrac{m \cdot v^2}{2}$ (massa vezes a velocidade ao quadrado, dividido por 2).
Para a conta ficar fácil, faça de conta que a velocidade da $\alpha$ vale 1. Então a do nêutron vale 4:
\begin{conta}
$E_\alpha = \dfrac{4 \times 1^2}{2} = \dfrac{4 \times 1}{2} = 2
\qquad\qquad
E_n = \dfrac{1 \times 4^2}{2} = \dfrac{1 \times 16}{2} = 8$
\end{conta}
Atenção ao quadrado: $4^2 = 4 \times 4 = 16$ (não é $4 \times 2$).
Comparando: 8 é \textbf{4 vezes} 2. A energia do nêutron é 4 vezes a da partícula $\alpha$.
(Com qualquer outra velocidade, a comparação daria o mesmo.)

\npasso{Reparta os 20\un{MeV}.}
A $\alpha$ fica com 1 parte e o nêutron com 4 partes. No total são $1 + 4 = 5$ partes iguais.
\begin{conta}
Cada parte: $20 \div 5 = 4\un{MeV}$
\qquad\qquad
Nêutron: $4 \times 4 = 16\un{MeV}$
\qquad\qquad
$\alpha$: $1 \times 4 = 4\un{MeV}$
\end{conta}
\confere{$16 + 4 = 20\un{MeV}$. Bateu com a energia total.}

\resposta{E}{O nêutron, por ser 4 vezes mais leve, fica com 4 das 5 partes da energia: 16\un{MeV}.}

\begin{outras}
\alt{A} 4,0\un{MeV} é a energia da partícula $\alpha$, não a do nêutron.\par
\alt{B} 10\un{MeV} seria dividir a energia meio a meio. Isso só valeria se as massas fossem iguais.\par
\altl{C} e \alt{D} Não saem de nenhuma divisão correta: a energia tem de ser repartida em 5 partes, 1 para a $\alpha$ e 4 para o nêutron.
\end{outras}

\begin{pegadinha}[Onde muita gente erra]
\begin{itemize}
  \item \textbf{Dividir a energia meio a meio} (10\un{MeV}). As massas são diferentes, então a divisão não é igual.
  \item \textbf{Achar que a partícula mais pesada leva mais energia.} É o contrário: a mais leve sai muito mais rápida, e a
        velocidade entra \emph{ao quadrado} na energia. Massas 4 e 1, energias 1 e 4.
\end{itemize}
\end{pegadinha}
